% F083｜06 总结：快记、快懂、快用
\clearpage
\begin{summarybox}
\textbf{\textcolor{CSummary}{一句话理解}}

\textbf{把 $x$ 换成 $-x$：指数会变倒数，根式中的一加一减也会互为倒数；绝对值和平方不变。看整个函数是“变相反数”还是“保持不变”，就能快速判定奇偶性。}

\textbf{\textcolor{CSummary}{快记：7 奇 + 5 偶}}

\textbf{奇函数（7 个）：}
\[
\begin{gathered}
 a^x-a^{-x},\quad \frac{a^x-1}{a^x+1},\quad
 \frac{a^x+1}{a^x-1},\\
 \ln\frac{kx-b}{kx+b},\quad \ln\frac{kx+b}{kx-b},\\
 \log_a(x+\sqrt{x^2+1}),\quad
 \log_a(\sqrt{x^2+1}-x).
\end{gathered}
\]
\textbf{偶函数（5 个）：}
\[
\begin{gathered}
 a^x+a^{-x},\quad a^{|x|},\quad a^{-x^2},\\
 \log_a|x|,\quad
 x\left(\frac12+\frac1{2^x-1}\right).
\end{gathered}
\]

\textbf{统一参数：}$a>0,a\ne1$；含 $k,b$ 的对数比值按 $k,b\ne0$；各式还要满足自己的定义域。
\end{summarybox}
\clearpage
\begin{summarybox}
\textbf{\textcolor{CSummary}{快懂：记住四个最短理由}}
\[
\begin{aligned}
 a^{-x}&=1/a^x,\\
 |{-x}|&=|x|,\quad(-x)^2=x^2,\\
 (\sqrt{x^2+1}+x)(\sqrt{x^2+1}-x)&=1,\\
 \log_a(1/t)&=-\log_a t\quad(t>0).
\end{aligned}
\]

\textbf{\textcolor{CSummary}{快用：解题三步}}
\begin{enumerate}[leftmargin=*,itemsep=.25em]
 \item \textbf{先看定义域：} $x$ 和 $-x$ 是否都能代？
 \item \textbf{再认结构：} 指数成对、分式倒数、平方差等于 $1$、$|x|$ 或 $x^2$、奇函数乘积。
 \item \textbf{最后一行证明：} 写出 $f(-x)=\pm f(x)$，并核对特殊点是否有定义。
\end{enumerate}

\textbf{\textcolor{CSummary}{三种证明怎么选？}}
\begin{description}[style=nextline,leftmargin=2.4em,labelsep=.5em,itemsep=.25em]
 \item[\textbf{门槛最低：}] 直接把 $x$ 换成 $-x$，适用于含 $|x|$、$x^2$ 的式子。
 \item[\textbf{最快：}] 先看是否“变倒数”，适用于指数分式、对数分式、根式对数。
 \item[\textbf{最能迁移：}] 先辨认每一部分的奇偶性，再用加减乘法规律处理复杂式子。
\end{description}

\textbf{\textcolor{CSummary}{与已有结论的联系}}

F015（定义域与奇偶性判定）是基础；F017（四则运算）解释“奇乘奇为偶”；F018（复合函数）解释绝对值与平方的情况；F047 关联 $a^{|x|}$；F009 可关联根式表达式。

\textbf{最后检查：}如果 $0$ 不在定义域，\textbf{不要套“奇函数在原点取 $0$”}；如果 $0$ 在定义域，才可用 $f(0)=0$ 帮助求参数。
\end{summarybox}
